\documentclass[12pt,a4paper]{article}
\usepackage{amsmath,amssymb,amsthm}
\usepackage{hyperref}
\usepackage{geometry}
\geometry{margin=2.5cm}
\usepackage{enumitem}

\newtheorem{theorem}{Theorem}[section]
\newtheorem{lemma}[theorem]{Lemma}
\newtheorem{corollary}[theorem]{Corollary}
\newtheorem{proposition}[theorem]{Proposition}
\theoremstyle{definition}
\newtheorem{definition}[theorem]{Definition}
\theoremstyle{remark}
\newtheorem{remark}[theorem]{Remark}

\title{Neutrino Mass Ordering from Recognition Science:\\
Normal Hierarchy Predicted by the $\varphi$-Ladder}

\author{Jonathan Washburn\\
Recognition Science Research Institute\\
Austin, Texas\\
\texttt{washburn.jonathan@gmail.com}}

\date{March 2026}

\begin{document}
\maketitle

\begin{abstract}
We prove that Recognition Science (RS) predicts normal mass
ordering for neutrinos: $m_1 < m_2 < m_3$.  The
$\varphi$-ladder assigns rung positions
$r_{\nu_1} = 0 < r_{\nu_2} = 11 < r_{\nu_3} = 19$.  Since
$\varphi > 1$, the mass formula $m_i = E_{\mathrm{coh}} \cdot
\varphi^{r_i}$ is strictly monotonic in $r$.  Inverted ordering
($m_3 < m_1$) would require $r_3 < r_1$, contradicting the
neutrino-sector rung gap structure derived from the $D = 3$ cube
geometry: the rung gaps $\Delta r_\nu(1) = 11$ and
$\Delta r_\nu(2) = 8$ are positive, ensuring $r_1 < r_2 < r_3$.
The prediction is falsifiable by JUNO, DUNE, and atmospheric
neutrino experiments.  All structural results are
machine-verified.
\end{abstract}

%% ============================================================
\section{Recognition Science Framework}
\label{sec:framework}
%% ============================================================

The normal-ordering proof depends on two RS structural results:
(i)~the $\varphi$-ladder uniqueness, which makes mass monotone in
rung index; and (ii)~the positive generation spacings derived from
the $D=3$ cube geometry.

\begin{definition}[RS Cost Functional]
For $x > 0$: $J(x) = \tfrac{1}{2}(x + x^{-1}) - 1$.
\end{definition}

\noindent\textbf{Axiom A1 (Normalization):} $J(1) = 0$.\\
\textbf{Axiom A2 (Recognition Composition Law, RCL):}
$J(xy)+J(x/y)=2J(x)J(y)+2J(x)+2J(y)$ for all $x,y>0$.\\
\textbf{Axiom A3 (Calibration):} $J''_{\log}(0) = 1$, where
$J_{\log}(t) := J(e^t)$.

\begin{theorem}[Cost Uniqueness]
\label{thm:unique}
The unique continuous solution to \textup{(A1)--(A3)} is
$J(x) = \tfrac{1}{2}(x + x^{-1}) - 1$.
\end{theorem}

\begin{proof}[Proof sketch]
Set $H(t) = J_{\log}(t)+1$.  Axiom A2 transforms into the d'Alembert
equation $H(s+t)+H(s-t)=2H(s)H(t)$.  By the Acz\'el classification
theorem~\cite{Aczel1966}, the unique continuous solution with $H(0)=1$
and $H''(0)=1$ is $H(t)=\cosh t$, giving
$J(x)=\tfrac{1}{2}(x+x^{-1})-1$.
\end{proof}

\noindent The $\varphi$-ladder places particles at rungs
$r \in \mathbb{Z}_{\geq 0}$ with mass $m = E_{\rm coh}\cdot\varphi^r$;
since $\varphi > 1$, mass is strictly increasing in rung.  The
$D = 3$ cube has three face-pair axes, forcing $N_{\rm gen} = 3$
generations with positive rung spacings between them.

%% ============================================================
\section{Introduction}
\label{sec:intro}
%% ============================================================

The neutrino mass ordering---whether $m_1 < m_2 < m_3$ (normal)
or $m_3 < m_1 < m_2$ (inverted)---is one of the most important
open questions in neutrino physics~\cite{PDG2024}.  Current data
from oscillation experiments show a mild preference for normal
ordering, but the determination is not yet
definitive~\cite{Capozzi2021}.

In the Standard Model, the mass ordering is a free parameter.
RS determines it uniquely.

%% ============================================================
\section{$\varphi$-Ladder Monotonicity}
\label{sec:mono}
%% ============================================================

\begin{definition}[$\varphi$-Ladder Mass Formula]
\begin{equation}
  m_i = E_{\mathrm{coh}} \cdot \varphi^{r_i},
  \qquad \varphi = \frac{1 + \sqrt{5}}{2} \approx 1.618.
\end{equation}
\end{definition}

\begin{lemma}[Strict Monotonicity]
\label{lem:mono}
If $r_i < r_j$, then $m_i < m_j$ (since $\varphi > 1$).
\end{lemma}

\begin{proof}
$m_j/m_i = \varphi^{r_j - r_i}$.  If $r_j > r_i$, then
$r_j - r_i > 0$, so $\varphi^{r_j - r_i} > 1$, hence
$m_j > m_i$.
\end{proof}

%% ============================================================
\section{Generation Spacing}
\label{sec:spacing}
%% ============================================================

\begin{definition}[Generation Spacings]
The $D = 3$ cube geometry determines neutrino-sector
rung gaps between successive generations:
\begin{align}
  \Delta r_\nu(1) &= 11
  \quad\text{(gap from }\nu_1\text{ to }\nu_2\text{)},\\
  \Delta r_\nu(2) &= 8
  \quad\text{(gap from }\nu_2\text{ to }\nu_3\text{)}.
\end{align}
Both gaps are positive, derived from the face-pair
structure of the $D = 3$ cube.
\end{definition}

\begin{remark}[Distinguishing neutrino and charged-lepton spacings]
The charged-lepton sector has rung gaps $\Delta r_\ell(1) = 11$
($e \to \mu$) and $\Delta r_\ell(2) = 17$ ($\mu \to \tau$).
The neutrino second gap $\Delta r_\nu(2) = 8$ differs from the
charged-lepton second gap $\Delta r_\ell(2) = 17$ due to the
absence of a gap correction for neutral particles ($Z_\nu = 0$).
\end{remark}

\begin{theorem}[Rung Ordering]
\label{thm:rung-order}
With gaps $\Delta r_\nu(1) = 11 > 0$ and $\Delta r_\nu(2) = 8 > 0$:
\begin{equation}
  r_{\nu_1} = 0, \quad
  r_{\nu_2} = r_{\nu_1} + \Delta r_\nu(1) = 11, \quad
  r_{\nu_3} = r_{\nu_2} + \Delta r_\nu(2) = 19.
\end{equation}
Therefore $r_{\nu_1} < r_{\nu_2} < r_{\nu_3}$.
\end{theorem}

%% ============================================================
\section{Normal Ordering Theorem}
\label{sec:normal}
%% ============================================================

\begin{theorem}[Normal Ordering]
\label{thm:normal}
\begin{equation}
  \boxed{m_{\nu_1} < m_{\nu_2} < m_{\nu_3}.}
\end{equation}
\end{theorem}

\begin{proof}
By Theorem~\ref{thm:rung-order}, $r_{\nu_1} < r_{\nu_2} <
r_{\nu_3}$.  By Lemma~\ref{lem:mono}, $m_{\nu_1} < m_{\nu_2} <
m_{\nu_3}$.
\end{proof}

%% ============================================================
\section{Exclusion of Inverted Ordering}
\label{sec:exclusion}
%% ============================================================

\begin{theorem}[Inverted Ordering Excluded]
\label{thm:inverted}
Inverted ordering ($m_3 < m_1$) is impossible in RS.
\end{theorem}

\begin{proof}
Inverted ordering requires $m_3 < m_1$, which by
Lemma~\ref{lem:mono} requires $r_3 < r_1$.  But
$r_3 = r_1 + \Delta r_\nu(1) + \Delta r_\nu(2)
= r_1 + 11 + 8 = r_1 + 19 > r_1$.  Contradiction.
\end{proof}

\begin{corollary}[Quasi-Degenerate Excluded]
The quasi-degenerate scenario ($m_1 \approx m_2 \approx m_3$) is
also excluded, since the rung separations are large
($\Delta r_\nu(1) = 11$, $\Delta r_\nu(2) = 8$), giving mass ratios
$m_2/m_1 = \varphi^{11} \gg 1$ and $m_3/m_2 = \varphi^{8} \gg 1$.
\end{corollary}

%% ============================================================
\section{Experimental Tests}
\label{sec:tests}
%% ============================================================

\begin{enumerate}
  \item \textbf{JUNO}~\cite{JUNO}: Reactor neutrino experiment
  sensitive to the mass ordering through $\Delta m_{31}^2$ vs.\
  $\Delta m_{32}^2$ separation.  Expected determination at
  $3\sigma$ within 6 years.

  \item \textbf{DUNE}~\cite{DUNE}: Long-baseline accelerator
  experiment.  Matter effects in $\nu_\mu \to \nu_e$ oscillations
  distinguish normal from inverted ordering at $5\sigma$.

  \item \textbf{Atmospheric neutrinos}:
  IceCube-Upgrade, KM3NeT/ORCA, and Hyper-K
  can determine the ordering via
  matter-enhanced oscillations.
\end{enumerate}

\begin{remark}
If any of these experiments determines inverted ordering,
Recognition Science is falsified.
\end{remark}

%% ============================================================
\section{Comparison}
\label{sec:comparison}
%% ============================================================

\begin{center}
\renewcommand{\arraystretch}{1.3}
\begin{tabular}{lcc}
\hline
\textbf{Framework} & \textbf{Ordering} &
\textbf{Prediction?} \\
\hline
Standard Model & Free parameter & No \\
See-saw mechanism & Either, model-dependent & Weak \\
$SO(10)$ GUT & Normal (typical) & Yes (model-dep.) \\
RS ($\varphi$-ladder) & Normal (theorem) & Yes (unique) \\
\hline
\end{tabular}
\end{center}

%% ============================================================
\section{Conclusion}
\label{sec:conclusion}
%% ============================================================

Recognition Science uniquely predicts normal mass ordering for
neutrinos.  The prediction follows from the strict monotonicity
of the $\varphi$-ladder ($\varphi > 1$) combined with the
positive neutrino-sector rung gaps
($\Delta r_\nu(1) = 11$, $\Delta r_\nu(2) = 8$)
derived from $D = 3$ cube geometry.  The prediction is sharp
and falsifiable: inverted ordering discovered by any experiment
would rule out RS.

%% ============================================================
\begin{thebibliography}{99}

\bibitem{RS_Axioms}
J.~Washburn, ``The Algebra of Reality: A Recognition Science Derivation
of Physical Law,'' \emph{Axioms} \textbf{15}(2), 90 (2026).
\texttt{doi:10.3390/axioms15020090}

\bibitem{Aczel1966}
J.~Acz\'el, \emph{Lectures on Functional Equations and Their
Applications} (Academic Press, 1966).

\bibitem{PDG2024}
S.~Navas \textit{et al.}\ (Particle Data Group),
``Review of Particle Physics,'' Phys.\ Rev.\ D \textbf{110}, 030001
(2024).

\bibitem{Capozzi2021}
F.~Capozzi, E.~Di Valentino, E.~Lisi, A.~Marrone,
A.~Melchiorri, and A.~Palazzo, ``Unfinished fabric of the three
neutrino paradigm,'' Phys.\ Rev.\ D \textbf{104}, 083031 (2021).

\bibitem{JUNO}
F.~An \textit{et al.}\ (JUNO), ``Neutrino Physics with JUNO,''
J.\ Phys.\ G \textbf{43}, 030401 (2016).

\bibitem{DUNE}
R.~Acciarri \textit{et al.}\ (DUNE), ``Long-Baseline Neutrino
Facility (LBNF) and Deep Underground Neutrino Experiment (DUNE),''
arXiv:1512.06148 (2015).

\end{thebibliography}

\end{document}
